% Bài làm buổi 18/11
\begin{problem}[9DS-HPT-002]
Từ phương trình thứ nhất: $x = 1 + 2y$.

Thế vào phương trình thứ hai: $3(1 + 2y) + y = 10$.

$3 + 7y = 10$, nên $7y = 7$, suy ra $y = 1$.

Thay $y = 1$ vào $x = 1 + 2y$: $x = 3$.

Thử lại: $3 - 2 \cdot 1 = 1$ và $3 \cdot 3 + 1 = 10$ (đúng).

Đáp số: $(x; y) = (3; 1)$
\end{problem}

\begin{problem}[9DS-HPT-005]
Gọi giá một quyển vở loại A là $x$ (đồng), loại B là $y$ (đồng), điều kiện $x > 0$, $y > 0$.

Theo đề bài ta có hệ: $3x + 2y = 54\,000$ và $5x + 2y = 69\,000$.

Trừ vế theo vế: $2x = 15\,000$, suy ra $x = 7\,500$; khi đó $2y = 54\,000 - 22\,500 = 31\,500$, suy ra $y = 15\,750$ (thoả mãn điều kiện).

Đáp số: Loại A: $7\,500$ đồng; loại B: $15\,750$ đồng
\end{problem}
